% !TEX program = xelatex
\documentclass[randomquote=true]{HiBeamer}


\usepackage[answer_shown]{ChoiceQuestion}
\usepackage{PhyUnit}

\title{高中数理综合讲义}
\author{授课教师}
\date{\today}



\begin{document}

\begin{frame}[plain]
  \titlepage
\end{frame}

% --- 在开头插入全局目录 ---
\begin{frame}{课程大纲}
  \tableofcontents
\end{frame}

\section{力学基础}

\begin{frame}{牛顿第二定律}
  \begin{defn}[牛顿第二定律]
    物体的加速度跟物体所受的合外力成正比，跟物体的质量成反比\footnote{这是动量定理在质量不变时的特例。}。
    \[ \vec{F}_{\text{合}} = m \vec{a} \]
  \end{defn}

  \begin{tip}[解题技巧]
    运用牛顿第二定律时，\hl{正交分解}是最核心的方法。
  \end{tip}
\end{frame}


\begin{frame}{123}
123  \xzanswer{B} \tkanswer{填空}

\fourchoices[answer=C]
{1}
{2}
{455676ghfh7}
{6}



\end{frame}


\begin{frame}{力学基础例题}
	% 默认继续上次的编号
	\begin{ExampleList}[0]
		\item 质量为 $m$ 的物体在水平面上受摩擦力减速，求加速度。
		
		\pause
		\item 关于牛顿第一定律，下列说法正确的是？
	\end{ExampleList}
\end{frame}

\begin{frame}{力学基础例题}
	% 默认继续上次的编号
	\begin{ExampleList}
		\item 质量为 $m$ 的物体在水平面上受摩擦力减速，求加速度。
		\pause 
		\item 关于牛顿第一定律，下列说法正确的是？
	\end{ExampleList}
\end{frame}



\begin{frame}{综合训练}
	\begin{ExerciseList}
		\item 斜面上的物体受力分析综合题。
		\item 合题11111
		\item  jkjk
	\end{ExerciseList}
	
\end{frame}

\begin{frame}{title}
	\begin{ExampleList}
		
		\item 质量为 $0.5 \Ukg$ 的0.5 kg 小球碰墙反弹，求冲量。质量为 $0.5\text{kg}$ 的小球碰墙反弹，求冲量。质量为 $0.5\Ukg$ 的小球碰墙反弹，求冲量。
		\begin{enumerate}
			\item 量为 $0.5\text{kg}$ 的小球量为 $0.5\text{kg}$ 的小球量为 $0.5\text{kg}$ 的小球量为 $0.5\text{kg}$ 的小球量为 $0.5\text{kg}$ 的小球量为 $0.5\text{kg}$ 的小球
			量为 $0.5\text{kg}$ 的小球
			\item 量为 $0.5\text{kg}$ 的小球
		\end{enumerate}
		\jdanswer{
			\begin{enumerate}
				\item 189
				\item  ioio
			\end{enumerate}
		}
		
	\item 质量为 $0.5\text{kg}$ 的小球碰墙反弹，求冲量。质量为 $0.5\text{kg}$ 的小球碰墙反弹，求冲量。质量为 $0.5\Ukg$ 的小球碰墙反弹，求冲量。
	\begin{enumerate}
		\item 量为 $0.5\text{kg}$ 的小球量为 $0.5\text{kg}$ 的小球量为 $0.5\text{kg}$ 的小球量为 $0.5\text{kg}$ 的小球量为 $0.5\text{kg}$ 的小球量为 $0.5\text{kg}$ 的小球
		量为 $0.5\text{kg}$ 的小球
		\item 量为 $0.5\text{kg}$ 的小球
	\end{enumerate}
	\jdanswer{
		\begin{enumerate}
			\item 189
			\item  ioio
		\end{enumerate}
	}
	\end{ExampleList}
	
\end{frame}

\begin{frame}{综合训练}
	\begin{ExampleList}[<+->]
		\item 斜面上的物体受力分析综合题。
		\item 合题
		\item  jkjk
	\end{ExampleList}
	\begin{ExerciseList}[<+->]
		\item 质量为 $0.5\text{kg}$ 的小球碰墙反弹，求冲量。
		\item 一颗子弹击中木块并留在其中，求共同速度。
	\end{ExerciseList}
\end{frame}


\begin{frame}{title}
	\begin{ExampleList}[<+->]
		\item 斜面上的物体受力分析综合题。
		\item 合题
		\item  jkjk
	\end{ExampleList}
\end{frame}

\begin{frame}{正交分解法步骤}
  运用正交分解法解题的标准流程：
  \begin{itemize}
  	\item 
  	90
  	\item 89
  \end{itemize}
  \begin{enumerate}
    \item 选取坐标系
    \begin{itemize}
      \item 沿运动方向（加速度方向）为 $x$ 轴
      \item 垂直运动方向为 $y$ 轴
    \end{itemize}
    \item 对各个力进行分解
    \item 列出分量方程
    \begin{enumerate}
      \item $x$ 轴方向：$F_x = ma_x$       \item $y$ 轴方向：$F_y = 0$ (通常平衡)
      
      \begin{enumerate}
      	\item 133
      	\item  test
      \end{enumerate}
      \begin{itemize}
      	\item 34
      	\item 56
      \end{itemize}
    \end{enumerate}
  \end{enumerate}
\end{frame}

\begin{frame}{斜面问题实战}
  \begin{eg}[斜面问题]
    质量为 $m$ 的物体在倾角为 $\theta$ 的斜面上匀速下滑，求动摩擦因数 $\mu$。
  \end{eg}

  受力分析步骤：
  \begin{itemize}[<+->]
    \item 确定研究对象
    \item 受力分析
    \begin{enumerate}[<+->]
      \item 重力 $G = mg$ 竖直向下
      \item 支持力 $N$ 垂直斜面向上 \tkanswer{8989} \tkanswer{8989} 
      
      \item 滑动摩擦力 $f$ 沿斜面向上
      \begin{enumerate}
      	\item 67
      	\item 45
      \end{enumerate}
    \end{enumerate}
    \item 建立坐标系并分解力
  \end{itemize}
\end{frame}

\begin{frame}{斜面问题拓展}
  \begin{supp}[拓展知识]
    如果斜面不是固定的，而是放在光滑水平面上，则需要考虑系统质心运动定理\footnote{此时两者构成的系统在水平方向动量守恒。}。
  \end{supp}
\end{frame}

\section{动量与冲量}

\begin{frame}{动量定理核心}
  \begin{imp}[核心考点]
    动量定理是矢量式，必须规定正方向！
    \[ \vec{I} = \Delta \vec{p} \]
  \end{imp}

  \begin{memo}[备课笔记]
    学生常犯错误：忘记 $\Delta \vec{v}$ 是末速度减初速度，尤其是反弹问题中速度方向改变导致的符号问题。
  \end{memo}
\end{frame}

\begin{frame}{冲量计算注意}
  \begin{warn}[易错警示]
    计算冲量时，\hlwarn{绝不能}直接用 $I = F t$，除非 $F$ 是恒力！对于变力，只能用动量定理求平均力。
  \end{warn}

  \begin{ex}[变力问题]
    质量为 $0.5\text{kg}$ 的小球从 $1.8\text{m}$ 高处自由落体，碰地后反弹高度为 $0.8\text{m}$，碰撞时间为 $0.01\text{s}$，求碰撞过程中的平均作用力。
  \end{ex}
\end{frame}

\begin{frame}{实验与数据分析}
  验证动量定理的实验数据：

  \vspace{0.3cm}
  \begin{table}
    \centering
    \begin{tabular}{cccc}
      \toprule
      次数 & 质量 $m/\text{kg}$ & 速度变化 $\Delta v/(\text{m/s})$ & 时间 $\Delta t/\text{s}$ \\
      \midrule
      1 & 0.50 & 6.0 & 0.12 \\
      2 & 0.50 & 8.0 & 0.16 \\
      3 & 1.00 & 6.0 & 0.24 \\
      \bottomrule
    \end{tabular}
    \caption{实验数据}
  \end{table}
\end{frame}

\begin{frame}{实验注意事项}
  \begin{itemize}
    \item 气垫导轨的调平\footnote{若未调平，重力分量会参与冲量计算，导致系统误差。遮光条越窄，速度测量越精确遮光条越窄，速度测量越精确遮光条越窄，速度测量越精确遮光条越窄，速度测量越精确遮光条越窄，速度测量越精确遮光条越窄，}
    \item 光电门间距的准确测量
    \item 遮光条宽度的补偿。遮光条越窄，速度测量越精确遮光条越窄，速度测量越精确遮光条越窄，速度测量越精确遮光条越窄，速度测量越精确遮光条越窄，速度测量越精确遮光条越窄，速度测量越精确遮光条越窄，速度测量越精确遮光条越窄，速度测量越精确遮光条越窄，速度测量越精确遮光条越窄，速度测量越精确
    \begin{itemize}
      \item 遮光条越窄，速度测量越精确遮光条越窄，速度测量越精确遮光条越窄，速度测量越精确遮光条越窄，速度测量越精确遮光条越窄，速度测量越精确遮光条越窄，速度测量越精确遮光条越窄，速度测量越精确遮光条越窄，速度测量越精确遮光条越窄，速度测量越精确遮光条越窄，速度测量越精确
      \item 但过窄会导致计时误差增大
    \end{itemize}
  \end{itemize}
\end{frame}

\begin{frame}{图像法应用}
  利用 $F-t$ 图像求冲量（面积法）：
To use a beamer-font, you can use the command.
利用 $F-t$ 图像求冲量（面积法）：
To use a beamer-font, you can use the command.
利用 $F-t$ 图像求冲量（面积法）：
To use a beamer-font, you can use the command.
  \begin{figure}
    \begin{tikzpicture}
      \draw[highend-grid] (-0.5,-0.5) grid (4,3);
      \draw[highend-axis] (-0.5,0) -- (4.5,0) node[right] {$t/\text{s}$};
      \draw[highend-axis] (0,-0.5) -- (0,3.5) node[above] {$F/\text{N}$};
      % 变力曲线
      \draw[highend-curve, domain=0:4, samples=50] plot (\x, {0.2*(\x)*(\x)});
      % 辅助线
      \draw[highend-aux] (4,0) -- (4,3.2);
      \node[below, font=\small] at (4,0) {$t_1$};
      \node[left, font=\small] at (0,3.2) {$F_1$};
      % 面积标注
      \node[text=AccentGold, font=\bfseries] at (2, 1) {阴影面积 $= I$};
      \fill[AccentGold, opacity=0.1] (0,0) -- plot[domain=0:4, samples=50] (\x, {0.2*(\x)*(\x)}) -- (4,0) -- cycle;
    \end{tikzpicture}
    \caption{图像求冲量}
  \end{figure}
\end{frame}


\begin{frame}{图表与公式编号测试}
  \begin{columns}[T]
    \begin{column}{0.5\linewidth}
    牛顿第二定律：遮光条越窄，速度测量越精确遮光条越窄，速度测量越精确遮光条越窄，
      % 测试图与子图
      \begin{figure}
        \centering
        \begin{subfigure}[b]{0.45\textwidth}
          \centering
          \begin{tikzpicture}
            \draw[highend-axis] (-0.2,0) -- (2.2,0);
            \draw[highend-curve, domain=0:2] plot (\x, {\x*\x});
          \end{tikzpicture}
          \caption{抛物线}
        \end{subfigure}
        \hfill
        \begin{subfigure}[b]{0.45\textwidth}
          \centering
          \begin{tikzpicture}
            \draw[highend-axis] (-0.2,0) -- (2.2,0);
            \draw[highend-curve, domain=0:1.5] plot (\x, {exp(\x)});
          \end{tikzpicture}
          \caption{指数函数}
        \end{subfigure}
        \caption{基本初等函数示例}
      \end{figure}
    \end{column}

    \begin{column}{0.45\linewidth}
      % 测试带编号的公式
      牛顿第二定律：遮光条越窄，速度测量越精确遮光条越窄，速度测量越精确遮光条越窄，速度测量越精确遮光条越窄
      \begin{equation}
        \vec{F} = m \vec{a}
      \end{equation}
      动能定理：
      \begin{equation}
        W_{\text{合}} = \Delta E_k
      \end{equation}
      度测量越精确遮光条越窄，速度测量越精确遮光条越窄，速度测量越精确
    \end{column}
  \end{columns}
\end{frame}

\end{document}
